\documentclass[leqno, a4paper, 12pt]{jsarticle}
\usepackage{amssymb}
\usepackage{amsthm}
\usepackage{amsmath}
\usepackage{comment}
\usepackage{url}

\usepackage{enumitem}
%\usepackage{showkeys}


%定理環境 thmはあまり使わずpropをなるべく使う。
%主張は基本的に証明の中で使い、そのため番号を付けない。
%注意環境を付けてみた。
\newtheorem{thm}{定理}
\newtheorem{prop}[thm]{命題}
\newtheorem{cor}[thm]{系}
\newtheorem{lem}[thm]{補題}
\newtheorem{df}[thm]{定義}
\newtheorem{exam}[thm]{例}
\newtheorem{fact}[thm]{事実}
\newtheorem*{claim}{主張}
\newtheorem*{cau}{注意}
\newtheorem*{rmk}{注意}
\renewcommand{\proofname}{証明}
%矢印たち。
%\toは\rightarrowと同じ。\getsは\leftarrowと同じ。同値は\iff
\newcommand{\To}{\Longrightarrow}
\newcommand{\Gets}{\LongLeftarrow}
%黒板太字
\newcommand{\nn}{\mathbb{N}}
\newcommand{\zz}{\mathbb{Z}}
\newcommand{\qq}{\mathbb{Q}}
\newcommand{\rr}{\mathbb{R}}
\newcommand{\cc}{\mathbb{C}}
%無限大
\newcommand{\mgn}{\infty}
%差集合は\setminusで統一する。等号の否定は\neq
%真部分集合は\subsetneqq　偏微分は\partial
%ディスプレイスタイル。\dfracでディスプレイ分数
\newcommand{\dpst}{\displaystyle}

\title{論文メモ
\large{\\--- Bing ---}
}
\author{ナンブ　キトラ}
\date{\today}
\begin{document}
\maketitle

本棚を整理したいので，
溜まっている論文のメモを書いて未来への手紙とすることにする．

R H Bing っていう，すげぇ超おもしれー数学をやる数学者がいたんだけどそのお仕事の一端のお話．

むかしむかし，どれくらい昔かというと
``$\mathbb{X}$''
が
``Twitter''
と呼ばれていて青い鳥もいてbotも普通に動いていた時代，
「トポロ沢さん」
という位相空間の問題を出すbotアカウントが
あったのだがその問題のうち私が解けなかった問題が，
位相空間$X$で$X\times \rr$が$\rr^{4}$
に同相なものは存在するか？という問題だった．
この問題を解くためには「ビング的なトポロジー」を勉強する必要が
ありこの問題に関する論文を
集めたのが今回の論文の束だと思う．



時系列順に並べるとこんな感じ．

1957： \cite{MR0092961}

この論文ではMoore decomposition 
とかいうupper semi-continuous decomposition 
(USCD 今ふうにいうと確か閉同値関係による分割)
を$\rr^{3}$で考えて，
その商空間が
$\rr^{3}$
に同相でないものを構成しているようだ．
ちなみに$\rr^{2}$だと
$\rr^{2}$に同相というのが
それこそMooreによって示されているようだ．

1957 ：
\cite{MR0092963}

この論文は上の論文を受けての
論文だと思う．

1957：
\cite{MR0092962}

この論文ではユークリッド空間の
いい感じの同相写像の構成をやっている．
この論文の結果からBing の商空間は
$\rr^{4}$に埋め込めるということがわかるようだ．

1958：\cite{MR0097034}

この論文は
\cite{MR0107228}の
アナウンスのような気がする．

1959：
\cite{MR0107228}

この論文では一番上の論文で構成した
商空間$B$について
$B\times \rr$が
$\rr^{4}$に同相であることを証明している．

1959：
\cite{MR0102817}

なんか色々やっている．
$\rr^{4}$の部分3次元多様体$M$で
$H_{1}(M)=0$だが
$\pi_{1}(M)\neq 0$
になるようなものを構成している．
また一番上の論文で構成された商空間$B$
が各点でユークリッド的でないにも関わらず
$3$次元のホモトピー多様体であることを証明している．

1961：
\cite{MR0123303}

論文の出だしが
``We are learning strange things about cartesian products''なのが面白い．
内容はちょっとよくわかんなかった．

%READ!
%myplain is the same to amsplain style. 
\nocite{*}
\bibliographystyle{myplaindoidoi}
\bibliography{Bingsp.bib}



\end{document}